% --- Archon physics formalization source begin ---
% archon:physics
% archon:covers PhyXMiniProblems/problem_phyx_mini_0816.lean
% archon:source-report reports/phyx_mini/problem_phyx_mini_0816.source.json

\chapter{Physics problem phyx\_mini\_0816}
\label{ch:PhyXMiniProblems_problem_phyx_mini_0816}

\paragraph{Problem source.}
\#\# Physical scenario

Figure shows an elastic collision of two pucks (masses \textbackslash{}( m\_\{\textbackslash{}text\{A\}\} = 0.500 \textbackslash{}text\{ kg\} \textbackslash{}) and \textbackslash{}( m\_\{\textbackslash{}text\{B\}\} = 0.300 \textbackslash{}text\{ kg\} \textbackslash{}) on a frictionless air-hockey table. Puck A has an initial velocity of \textbackslash{}( 4.00 \textbackslash{}text\{ m/s\} \textbackslash{}) in the positive \textbackslash{}( x \textbackslash{})-direction and a final velocity of \textbackslash{}( 2.00 \textbackslash{}text\{ m/s\} \textbackslash{}) in an unknown direction \textbackslash{}( \textbackslash{}alpha \textbackslash{}). Puck B is initially at rest.

\#\# Question

Find the final speed \textbackslash{}( v\_\{\textbackslash{}text\{B2\}\} \textbackslash{}) of puck B.

\#\# Answer choices

- A: A: 3.81m/s

- B: B: 4.47m/s

- C: C: 2.96m/s

- D: D: 3.62m/s

\#\# Figure caption (auxiliary; use the image as primary evidence)

The image consists of two parts showing a collision between two objects, A and B, on the x-y plane.

**Before Collision:**
- Object A is moving along the x-axis with velocity \textbackslash{}( v\_\{A1\} = 4.00 \textbackslash{}, \textbackslash{}text\{m/s\} \textbackslash{}).
- Object B is stationary.
- The mass of A is \textbackslash{}( m\_A = 0.500 \textbackslash{}, \textbackslash{}text\{kg\} \textbackslash{}), and the mass of B is \textbackslash{}( m\_B = 0.300 \textbackslash{}, \textbackslash{}text\{kg\} \textbackslash{}).

**After Collision:**
- Object A moves with a velocity \textbackslash{}( v\_\{A2\} = 2.00 \textbackslash{}, \textbackslash{}text\{m/s\} \textbackslash{}) at an angle \textbackslash{}( \textbackslash{}alpha \textbackslash{}) from the x-axis.
- Object B moves with a velocity \textbackslash{}( v\_\{B2\} \textbackslash{}) at an angle \textbackslash{}( \textbackslash{}beta \textbackslash{}) from the x-axis.
- Arrows represent the direction of motion.

Both parts are labeled "BEFORE" and "AFTER" to indicate the sequence of events. The scene reflects a typical physics problem involving conservation of momentum in two-dimensional collisions.

\#\# Dataset metadata

Momentum and Collisions; Physical Model Grounding Reasoning, Multi-Formula Reasoning

\paragraph{Recorded answer/context.}
B: B: 4.47m/s

\paragraph{Figure/image path.}
/root/proposal\_for\_physic/science-mango/phyx\_mini\_run/phyx\_data/test\_image/816.png

\paragraph{Formalization target.}
create a compiling Lean file with sorry bodies at `PhyXMiniProblems/problem\_phyx\_mini\_0816.lean`. The Lean declarations must preserve the physical quantities, dimensions or dimensional roles, figure labels, governing-law hypotheses, and final relation expressed by this problem.
Use Mathlib/Physlib names found through LeanExplore where available. If a physics API is missing, introduce faithful local abstractions rather than scalar placeholder aliases.

\begin{theorem}[Physics formalization target]
\label{thm:physics:phyx_mini_0816:target}
\lean{PhyXMiniProblems.ProblemPhyXMini0816.elasticPuckCollision_finalSpeed}
\uses{def:physics:phyx-mini-0816:phyxminiproblems-problemphyxmini0816-speedmagnitudeinmeterspersecond, def:physics:phyx-mini-0816:phyxminiproblems-problemphyxmini0816-elasticpuckcollisionsetup, def:physics:phyx-mini-0816:phyxminiproblems-problemphyxmini0816-matcheselasticairhockeyscenario, def:physics:phyx-mini-0816:phyxminiproblems-problemphyxmini0816-matchespuckcollisionproblemdata, def:physics:phyx-mini-0816:phyxminiproblems-problemphyxmini0816-hasphysicalpuckmasses, def:physics:phyx-mini-0816:phyxminiproblems-problemphyxmini0816-matchessuppliedpuckcollisionfigure, def:physics:phyx-mini-0816:phyxminiproblems-problemphyxmini0816-hasvelocitymagnitudekinematics, def:physics:phyx-mini-0816:phyxminiproblems-problemphyxmini0816-matchesfigurevelocitygeometry, def:physics:phyx-mini-0816:phyxminiproblems-problemphyxmini0816-satisfieselasticcollisionconservationlaws, def:physics:phyx-mini-0816:phyxminiproblems-problemphyxmini0816-matchesdisplayedspeed, def:physics:phyx-mini-0816:phyxminiproblems-problemphyxmini0816-isuniqueclosestspeedchoice}
The assigned autoformalize agent should translate this physics problem into Lean declarations in the covered file, with theorem and lemma proof bodies written as `by sorry`.
\end{theorem}
\begin{proof}
This is an autoformalization task, not a proof task. Produce statements that can later be proved by the physics prover without weakening the physical model.
\end{proof}

% --- Archon declaration topology begin ---
\section{Declaration topology}
\begin{definition}[velocity Dimension]
  \label{def:physics:phyx-mini-0816:phyxminiproblems-problemphyxmini0816-velocitydimension}
  \lean{PhyXMiniProblems.ProblemPhyXMini0816.velocityDimension}
  The physical dimension of velocity, L T⁻¹
\end{definition}
\begin{proof}
  This is the declared model definition of velocity Dimension.
\end{proof}
\begin{definition}[Mass Quantity]
  \label{def:physics:phyx-mini-0816:phyxminiproblems-problemphyxmini0816-massquantity}
  \lean{PhyXMiniProblems.ProblemPhyXMini0816.MassQuantity}
  A nonnegative, unit-independent physical mass
\end{definition}
\begin{proof}
  This is the declared model definition of Mass Quantity.
\end{proof}
\begin{definition}[Speed Magnitude Quantity]
  \label{def:physics:phyx-mini-0816:phyxminiproblems-problemphyxmini0816-speedmagnitudequantity}
  \lean{PhyXMiniProblems.ProblemPhyXMini0816.SpeedMagnitudeQuantity}
  A nonnegative, unit-independent physical speed magnitude
\end{definition}
\begin{proof}
  This is the declared model definition of Speed Magnitude Quantity.
\end{proof}
\begin{definition}[Planar Velocity Quantity]
  \label{def:physics:phyx-mini-0816:phyxminiproblems-problemphyxmini0816-planarvelocityquantity}
  \lean{PhyXMiniProblems.ProblemPhyXMini0816.PlanarVelocityQuantity}
  \uses{def:physics:phyx-mini-0816:phyxminiproblems-problemphyxmini0816-velocitydimension}
  A signed, unit-independent velocity vector in the table plane
\end{definition}
\begin{proof}
  This is the declared model definition of Planar Velocity Quantity.
\end{proof}
\begin{definition}[x Axis]
  \label{def:physics:phyx-mini-0816:phyxminiproblems-problemphyxmini0816-xaxis}
  \lean{PhyXMiniProblems.ProblemPhyXMini0816.xAxis}
  Coordinate 0, the positive horizontal axis pointing right
\end{definition}
\begin{proof}
  This is the declared model definition of x Axis.
\end{proof}
\begin{definition}[y Axis]
  \label{def:physics:phyx-mini-0816:phyxminiproblems-problemphyxmini0816-yaxis}
  \lean{PhyXMiniProblems.ProblemPhyXMini0816.yAxis}
  Coordinate 1, the positive vertical axis pointing up
\end{definition}
\begin{proof}
  This is the declared model definition of y Axis.
\end{proof}
\begin{definition}[mass Readout]
  \label{def:physics:phyx-mini-0816:phyxminiproblems-problemphyxmini0816-massreadout}
  \lean{PhyXMiniProblems.ProblemPhyXMini0816.massReadout}
  \uses{def:physics:phyx-mini-0816:phyxminiproblems-problemphyxmini0816-massquantity}
  Read a physical mass in a selected mass unit
\end{definition}
\begin{proof}
  This is the declared model definition of mass Readout.
\end{proof}
\begin{definition}[speed Magnitude Readout]
  \label{def:physics:phyx-mini-0816:phyxminiproblems-problemphyxmini0816-speedmagnitudereadout}
  \lean{PhyXMiniProblems.ProblemPhyXMini0816.speedMagnitudeReadout}
  \uses{def:physics:phyx-mini-0816:phyxminiproblems-problemphyxmini0816-speedmagnitudequantity}
  Read a speed magnitude in coherent selected length and time units
\end{definition}
\begin{proof}
  This is the declared model definition of speed Magnitude Readout.
\end{proof}
\begin{definition}[planar Velocity Readout]
  \label{def:physics:phyx-mini-0816:phyxminiproblems-problemphyxmini0816-planarvelocityreadout}
  \lean{PhyXMiniProblems.ProblemPhyXMini0816.planarVelocityReadout}
  \uses{def:physics:phyx-mini-0816:phyxminiproblems-problemphyxmini0816-planarvelocityquantity}
  Read a planar velocity in coherent selected length and time units
\end{definition}
\begin{proof}
  This is the declared model definition of planar Velocity Readout.
\end{proof}
\begin{definition}[mass In Kilograms]
  \label{def:physics:phyx-mini-0816:phyxminiproblems-problemphyxmini0816-massinkilograms}
  \lean{PhyXMiniProblems.ProblemPhyXMini0816.massInKilograms}
  \uses{def:physics:phyx-mini-0816:phyxminiproblems-problemphyxmini0816-massquantity, def:physics:phyx-mini-0816:phyxminiproblems-problemphyxmini0816-massreadout}
  Kilogram readout of a physical mass
\end{definition}
\begin{proof}
  This is the declared model definition of mass In Kilograms.
\end{proof}
\begin{definition}[speed Magnitude In Meters Per Second]
  \label{def:physics:phyx-mini-0816:phyxminiproblems-problemphyxmini0816-speedmagnitudeinmeterspersecond}
  \lean{PhyXMiniProblems.ProblemPhyXMini0816.speedMagnitudeInMetersPerSecond}
  \uses{def:physics:phyx-mini-0816:phyxminiproblems-problemphyxmini0816-speedmagnitudequantity, def:physics:phyx-mini-0816:phyxminiproblems-problemphyxmini0816-speedmagnitudereadout}
  Metre-per-second readout of a physical speed magnitude
\end{definition}
\begin{proof}
  This is the declared model definition of speed Magnitude In Meters Per Second.
\end{proof}
\begin{definition}[planar Velocity In Meters Per Second]
  \label{def:physics:phyx-mini-0816:phyxminiproblems-problemphyxmini0816-planarvelocityinmeterspersecond}
  \lean{PhyXMiniProblems.ProblemPhyXMini0816.planarVelocityInMetersPerSecond}
  \uses{def:physics:phyx-mini-0816:phyxminiproblems-problemphyxmini0816-planarvelocityquantity, def:physics:phyx-mini-0816:phyxminiproblems-problemphyxmini0816-planarvelocityreadout}
  Metre-per-second readout of a physical planar velocity
\end{definition}
\begin{proof}
  This is the declared model definition of planar Velocity In Meters Per Second.
\end{proof}
\begin{definition}[Puck Label]
  \label{def:physics:phyx-mini-0816:phyxminiproblems-problemphyxmini0816-pucklabel}
  \lean{PhyXMiniProblems.ProblemPhyXMini0816.PuckLabel}
  The two colliding pucks labelled in the problem
\end{definition}
\begin{proof}
  This is the declared model definition of Puck Label.
\end{proof}
\begin{definition}[Collision Stage]
  \label{def:physics:phyx-mini-0816:phyxminiproblems-problemphyxmini0816-collisionstage}
  \lean{PhyXMiniProblems.ProblemPhyXMini0816.CollisionStage}
  The two panels in the supplied before/after collision diagram
\end{definition}
\begin{proof}
  This is the declared model definition of Collision Stage.
\end{proof}
\begin{definition}[Diagram Axis]
  \label{def:physics:phyx-mini-0816:phyxminiproblems-problemphyxmini0816-diagramaxis}
  \lean{PhyXMiniProblems.ProblemPhyXMini0816.DiagramAxis}
  The coordinate axes drawn through the origin O in both panels
\end{definition}
\begin{proof}
  This is the declared model definition of Diagram Axis.
\end{proof}
\begin{definition}[Velocity Arrow Direction]
  \label{def:physics:phyx-mini-0816:phyxminiproblems-problemphyxmini0816-velocityarrowdirection}
  \lean{PhyXMiniProblems.ProblemPhyXMini0816.VelocityArrowDirection}
  Qualitative directions of the three velocity arrows in the raster
\end{definition}
\begin{proof}
  This is the declared model definition of Velocity Arrow Direction.
\end{proof}
\begin{definition}[Outgoing Angle Label]
  \label{def:physics:phyx-mini-0816:phyxminiproblems-problemphyxmini0816-outgoinganglelabel}
  \lean{PhyXMiniProblems.ProblemPhyXMini0816.OutgoingAngleLabel}
  The two outgoing angle labels printed in the after-collision panel
\end{definition}
\begin{proof}
  This is the declared model definition of Outgoing Angle Label.
\end{proof}
\begin{definition}[puck For Outgoing Angle]
  \label{def:physics:phyx-mini-0816:phyxminiproblems-problemphyxmini0816-puckforoutgoingangle}
  \lean{PhyXMiniProblems.ProblemPhyXMini0816.puckForOutgoingAngle}
  \uses{def:physics:phyx-mini-0816:phyxminiproblems-problemphyxmini0816-pucklabel, def:physics:phyx-mini-0816:phyxminiproblems-problemphyxmini0816-outgoinganglelabel}
  The outgoing puck whose direction is named by each printed angle
\end{definition}
\begin{proof}
  This is the declared model definition of puck For Outgoing Angle.
\end{proof}
\begin{definition}[direction For Outgoing Angle]
  \label{def:physics:phyx-mini-0816:phyxminiproblems-problemphyxmini0816-directionforoutgoingangle}
  \lean{PhyXMiniProblems.ProblemPhyXMini0816.directionForOutgoingAngle}
  \uses{def:physics:phyx-mini-0816:phyxminiproblems-problemphyxmini0816-velocityarrowdirection, def:physics:phyx-mini-0816:phyxminiproblems-problemphyxmini0816-outgoinganglelabel}
  The qualitative outgoing direction associated with each angle label
\end{definition}
\begin{proof}
  This is the declared model definition of direction For Outgoing Angle.
\end{proof}
\begin{definition}[Puck Collision Figure]
  \label{def:physics:phyx-mini-0816:phyxminiproblems-problemphyxmini0816-puckcollisionfigure}
  \lean{PhyXMiniProblems.ProblemPhyXMini0816.PuckCollisionFigure}
  \uses{def:physics:phyx-mini-0816:phyxminiproblems-problemphyxmini0816-pucklabel, def:physics:phyx-mini-0816:phyxminiproblems-problemphyxmini0816-collisionstage, def:physics:phyx-mini-0816:phyxminiproblems-problemphyxmini0816-diagramaxis, def:physics:phyx-mini-0816:phyxminiproblems-problemphyxmini0816-velocityarrowdirection, def:physics:phyx-mini-0816:phyxminiproblems-problemphyxmini0816-outgoinganglelabel}
  Literal and schematic content transcribed from the primary image. The displayedSpeedMetersPerSecond field is optional because the raster prints no numerical value next to B's outgoing velocity arrow
\end{definition}
\begin{proof}
  This is the declared model definition of Puck Collision Figure.
\end{proof}
\begin{definition}[Table Surface Model]
  \label{def:physics:phyx-mini-0816:phyxminiproblems-problemphyxmini0816-tablesurfacemodel}
  \lean{PhyXMiniProblems.ProblemPhyXMini0816.TableSurfaceModel}
  The idealized table model stated in the prose
\end{definition}
\begin{proof}
  This is the declared model definition of Table Surface Model.
\end{proof}
\begin{definition}[Collision Model]
  \label{def:physics:phyx-mini-0816:phyxminiproblems-problemphyxmini0816-collisionmodel}
  \lean{PhyXMiniProblems.ProblemPhyXMini0816.CollisionModel}
  The collision model stated in the prose
\end{definition}
\begin{proof}
  This is the declared model definition of Collision Model.
\end{proof}
\begin{definition}[Elastic Puck Collision Setup]
  \label{def:physics:phyx-mini-0816:phyxminiproblems-problemphyxmini0816-elasticpuckcollisionsetup}
  \lean{PhyXMiniProblems.ProblemPhyXMini0816.ElasticPuckCollisionSetup}
  \uses{def:physics:phyx-mini-0816:phyxminiproblems-problemphyxmini0816-massquantity, def:physics:phyx-mini-0816:phyxminiproblems-problemphyxmini0816-speedmagnitudequantity, def:physics:phyx-mini-0816:phyxminiproblems-problemphyxmini0816-planarvelocityquantity, def:physics:phyx-mini-0816:phyxminiproblems-problemphyxmini0816-pucklabel, def:physics:phyx-mini-0816:phyxminiproblems-problemphyxmini0816-collisionstage, def:physics:phyx-mini-0816:phyxminiproblems-problemphyxmini0816-puckcollisionfigure, def:physics:phyx-mini-0816:phyxminiproblems-problemphyxmini0816-tablesurfacemodel, def:physics:phyx-mini-0816:phyxminiproblems-problemphyxmini0816-collisionmodel}
  Independent physical fields of the experiment. In particular, B's outgoing speed is an independent dimensionful field and is not defined from a displayed answer value
\end{definition}
\begin{proof}
  This is the declared model definition of Elastic Puck Collision Setup.
\end{proof}
\begin{definition}[Matches Elastic Air Hockey Scenario]
  \label{def:physics:phyx-mini-0816:phyxminiproblems-problemphyxmini0816-matcheselasticairhockeyscenario}
  \lean{PhyXMiniProblems.ProblemPhyXMini0816.MatchesElasticAirHockeyScenario}
  \uses{def:physics:phyx-mini-0816:phyxminiproblems-problemphyxmini0816-elasticpuckcollisionsetup}
  The qualitative physical scenario supplied in the problem statement
\end{definition}
\begin{proof}
  This is the declared model definition of Matches Elastic Air Hockey Scenario.
\end{proof}
\begin{definition}[Matches Puck Collision Problem Data]
  \label{def:physics:phyx-mini-0816:phyxminiproblems-problemphyxmini0816-matchespuckcollisionproblemdata}
  \lean{PhyXMiniProblems.ProblemPhyXMini0816.MatchesPuckCollisionProblemData}
  \uses{def:physics:phyx-mini-0816:phyxminiproblems-problemphyxmini0816-massinkilograms, def:physics:phyx-mini-0816:phyxminiproblems-problemphyxmini0816-speedmagnitudeinmeterspersecond, def:physics:phyx-mini-0816:phyxminiproblems-problemphyxmini0816-elasticpuckcollisionsetup}
  Numerical physical readouts stated in the prose and shown in the image. No outgoing speed for puck B, exact result, rounded answer, or answer label occurs in this record
\end{definition}
\begin{proof}
  This is the declared model definition of Matches Puck Collision Problem Data.
\end{proof}
\begin{definition}[Has Physical Puck Masses]
  \label{def:physics:phyx-mini-0816:phyxminiproblems-problemphyxmini0816-hasphysicalpuckmasses}
  \lean{PhyXMiniProblems.ProblemPhyXMini0816.HasPhysicalPuckMasses}
  \uses{def:physics:phyx-mini-0816:phyxminiproblems-problemphyxmini0816-massinkilograms, def:physics:phyx-mini-0816:phyxminiproblems-problemphyxmini0816-elasticpuckcollisionsetup}
  Positive, nondegenerate physical masses for the two pucks
\end{definition}
\begin{proof}
  This is the declared model definition of Has Physical Puck Masses.
\end{proof}
\begin{definition}[Matches Supplied Puck Collision Figure]
  \label{def:physics:phyx-mini-0816:phyxminiproblems-problemphyxmini0816-matchessuppliedpuckcollisionfigure}
  \lean{PhyXMiniProblems.ProblemPhyXMini0816.MatchesSuppliedPuckCollisionFigure}
  \uses{def:physics:phyx-mini-0816:phyxminiproblems-problemphyxmini0816-elasticpuckcollisionsetup}
  Primary-image evidence from 816.png. The image supplies the mass and two A-speed labels, identifies B as initially at rest, and shows B's final arrow without a numerical speed label. Thus this structure cannot contain the requested answer
\end{definition}
\begin{proof}
  This is the declared model definition of Matches Supplied Puck Collision Figure.
\end{proof}
\begin{definition}[Has Velocity Magnitude Kinematics]
  \label{def:physics:phyx-mini-0816:phyxminiproblems-problemphyxmini0816-hasvelocitymagnitudekinematics}
  \lean{PhyXMiniProblems.ProblemPhyXMini0816.HasVelocityMagnitudeKinematics}
  \uses{def:physics:phyx-mini-0816:phyxminiproblems-problemphyxmini0816-speedmagnitudereadout, def:physics:phyx-mini-0816:phyxminiproblems-problemphyxmini0816-planarvelocityreadout, def:physics:phyx-mini-0816:phyxminiproblems-problemphyxmini0816-pucklabel, def:physics:phyx-mini-0816:phyxminiproblems-problemphyxmini0816-collisionstage, def:physics:phyx-mini-0816:phyxminiproblems-problemphyxmini0816-elasticpuckcollisionsetup}
  Every scalar speed field is the Euclidean magnitude of its associated planar velocity vector. This is a kinematic relation, not a numerical answer for B
\end{definition}
\begin{proof}
  This is the declared model definition of Has Velocity Magnitude Kinematics.
\end{proof}
\begin{definition}[Matches Figure Velocity Geometry]
  \label{def:physics:phyx-mini-0816:phyxminiproblems-problemphyxmini0816-matchesfigurevelocitygeometry}
  \lean{PhyXMiniProblems.ProblemPhyXMini0816.MatchesFigureVelocityGeometry}
  \uses{def:physics:phyx-mini-0816:phyxminiproblems-problemphyxmini0816-xaxis, def:physics:phyx-mini-0816:phyxminiproblems-problemphyxmini0816-yaxis, def:physics:phyx-mini-0816:phyxminiproblems-problemphyxmini0816-speedmagnitudereadout, def:physics:phyx-mini-0816:phyxminiproblems-problemphyxmini0816-planarvelocityreadout, def:physics:phyx-mini-0816:phyxminiproblems-problemphyxmini0816-elasticpuckcollisionsetup}
  Cartesian realization of the directions shown in the image. alpha is measured counterclockwise above positive x; beta is represented as a positive magnitude measured clockwise below positive x, hence B's negative y component
\end{definition}
\begin{proof}
  This is the declared model definition of Matches Figure Velocity Geometry.
\end{proof}
\begin{definition}[Satisfies Elastic Collision Conservation Laws]
  \label{def:physics:phyx-mini-0816:phyxminiproblems-problemphyxmini0816-satisfieselasticcollisionconservationlaws}
  \lean{PhyXMiniProblems.ProblemPhyXMini0816.SatisfiesElasticCollisionConservationLaws}
  \uses{def:physics:phyx-mini-0816:phyxminiproblems-problemphyxmini0816-massreadout, def:physics:phyx-mini-0816:phyxminiproblems-problemphyxmini0816-speedmagnitudereadout, def:physics:phyx-mini-0816:phyxminiproblems-problemphyxmini0816-planarvelocityreadout, def:physics:phyx-mini-0816:phyxminiproblems-problemphyxmini0816-elasticpuckcollisionsetup}
  The governing laws for an isolated elastic two-puck collision. Momentum is a two-dimensional vector equation. Kinetic energy uses the classical (1/2) m v² expression. Both laws hold in arbitrary coherent base units and constrain, but do not state, B's requested final speed
\end{definition}
\begin{proof}
  This is the declared model definition of Satisfies Elastic Collision Conservation Laws.
\end{proof}
\begin{definition}[Answer Choice]
  \label{def:physics:phyx-mini-0816:phyxminiproblems-problemphyxmini0816-answerchoice}
  \lean{PhyXMiniProblems.ProblemPhyXMini0816.AnswerChoice}
  The four speed-answer labels displayed with the problem
\end{definition}
\begin{proof}
  This is the declared model definition of Answer Choice.
\end{proof}
\begin{definition}[speed Meters Per Second]
  \label{def:physics:phyx-mini-0816:phyxminiproblems-problemphyxmini0816-answerchoice-speedmeterspersecond}
  \lean{PhyXMiniProblems.ProblemPhyXMini0816.AnswerChoice.speedMetersPerSecond}
  \uses{def:physics:phyx-mini-0816:phyxminiproblems-problemphyxmini0816-answerchoice}
  Speed in m/s printed beside each displayed answer choice
\end{definition}
\begin{proof}
  This is the declared model definition of speed Meters Per Second.
\end{proof}
\begin{definition}[recorded Dataset Answer]
  \label{def:physics:phyx-mini-0816:phyxminiproblems-problemphyxmini0816-recordeddatasetanswer}
  \lean{PhyXMiniProblems.ProblemPhyXMini0816.recordedDatasetAnswer}
  \uses{def:physics:phyx-mini-0816:phyxminiproblems-problemphyxmini0816-answerchoice}
  Dataset answer metadata, deliberately not used as a theorem premise
\end{definition}
\begin{proof}
  This is the declared model definition of recorded Dataset Answer.
\end{proof}
\begin{definition}[answer Choice Error Meters Per Second]
  \label{def:physics:phyx-mini-0816:phyxminiproblems-problemphyxmini0816-answerchoiceerrormeterspersecond}
  \lean{PhyXMiniProblems.ProblemPhyXMini0816.answerChoiceErrorMetersPerSecond}
  \uses{def:physics:phyx-mini-0816:phyxminiproblems-problemphyxmini0816-speedmagnitudeinmeterspersecond, def:physics:phyx-mini-0816:phyxminiproblems-problemphyxmini0816-elasticpuckcollisionsetup, def:physics:phyx-mini-0816:phyxminiproblems-problemphyxmini0816-answerchoice}
  Absolute discrepancy between B's physical final speed and one choice
\end{definition}
\begin{proof}
  This is the declared model definition of answer Choice Error Meters Per Second.
\end{proof}
\begin{definition}[Matches Displayed Speed]
  \label{def:physics:phyx-mini-0816:phyxminiproblems-problemphyxmini0816-matchesdisplayedspeed}
  \lean{PhyXMiniProblems.ProblemPhyXMini0816.MatchesDisplayedSpeed}
  \uses{def:physics:phyx-mini-0816:phyxminiproblems-problemphyxmini0816-elasticpuckcollisionsetup, def:physics:phyx-mini-0816:phyxminiproblems-problemphyxmini0816-answerchoice, def:physics:phyx-mini-0816:phyxminiproblems-problemphyxmini0816-answerchoiceerrormeterspersecond}
  Agreement with a speed displayed to the nearest 0.01 m/s
\end{definition}
\begin{proof}
  This is the declared model definition of Matches Displayed Speed.
\end{proof}
\begin{definition}[Is Unique Closest Speed Choice]
  \label{def:physics:phyx-mini-0816:phyxminiproblems-problemphyxmini0816-isuniqueclosestspeedchoice}
  \lean{PhyXMiniProblems.ProblemPhyXMini0816.IsUniqueClosestSpeedChoice}
  \uses{def:physics:phyx-mini-0816:phyxminiproblems-problemphyxmini0816-elasticpuckcollisionsetup, def:physics:phyx-mini-0816:phyxminiproblems-problemphyxmini0816-answerchoice, def:physics:phyx-mini-0816:phyxminiproblems-problemphyxmini0816-answerchoiceerrormeterspersecond}
  A displayed speed is strictly closer than every alternative
\end{definition}
\begin{proof}
  This is the declared model definition of Is Unique Closest Speed Choice.
\end{proof}
\begin{lemma}[puck BFinal Speed Squared]
  \label{lem:physics:phyx-mini-0816:phyxminiproblems-problemphyxmini0816-puckbfinalspeedsquared}
  \lean{PhyXMiniProblems.ProblemPhyXMini0816.puckBFinalSpeedSquared}
  \uses{def:physics:phyx-mini-0816:phyxminiproblems-problemphyxmini0816-speedmagnitudeinmeterspersecond, def:physics:phyx-mini-0816:phyxminiproblems-problemphyxmini0816-elasticpuckcollisionsetup, def:physics:phyx-mini-0816:phyxminiproblems-problemphyxmini0816-matchespuckcollisionproblemdata, def:physics:phyx-mini-0816:phyxminiproblems-problemphyxmini0816-satisfieselasticcollisionconservationlaws}
  Energy conservation and the stated readouts determine the square of B's final speed to be 20 (m/s)²
\end{lemma}
\begin{proof}
  The conclusion follows from the governing relations and figure readouts named in the statement of puck BFinal Speed Squared.
\end{proof}
\begin{lemma}[puck BFinal Speed Exact]
  \label{lem:physics:phyx-mini-0816:phyxminiproblems-problemphyxmini0816-puckbfinalspeedexact}
  \lean{PhyXMiniProblems.ProblemPhyXMini0816.puckBFinalSpeedExact}
  \uses{def:physics:phyx-mini-0816:phyxminiproblems-problemphyxmini0816-speedmagnitudeinmeterspersecond, def:physics:phyx-mini-0816:phyxminiproblems-problemphyxmini0816-elasticpuckcollisionsetup, def:physics:phyx-mini-0816:phyxminiproblems-problemphyxmini0816-matchespuckcollisionproblemdata, def:physics:phyx-mini-0816:phyxminiproblems-problemphyxmini0816-satisfieselasticcollisionconservationlaws}
  B's exact final speed is sqrt 20 m/s
\end{lemma}
\begin{proof}
  The conclusion follows from the governing relations and figure readouts named in the statement of puck BFinal Speed Exact.
\end{proof}
% --- Archon declaration topology end ---
% --- Archon physics formalization source end ---
